\documentclass{article}
\usepackage{PRIMEarxiv} % Core package for the PrimeAI Template
\usepackage{amsmath, amssymb, amsthm, bm, dcolumn} % Add amsthm here for the proof environment
\usepackage[numbers,sort&compress]{natbib} % Natbib for citations
\usepackage{graphicx} % For high-quality images
\usepackage{multicol} % Multi-column figures/tables
\usepackage{hyperref} % Hyperlinks
\usepackage{listings} % Code listings
\usepackage{authblk} % For structured affiliations
% Define theorem style (optional)
\theoremstyle{plain}
\newtheorem{theorem}{Theorem}
\renewcommand{\qedsymbol}{}

% Custom commands for primes and qubit encoding
\newcommand{\primeQubit}[1]{|p_{#1}\rangle}
\newcommand{\primeGate}[1]{U_{p_{#1}}}

\usepackage{wrapfig}
\usepackage[pscoord]{eso-pic}
\usepackage[fulladjust]{marginnote}
\reversemarginpar

% Typesetting improvements without footnote patching
\usepackage[protrusion=true, expansion=true, tracking=false]{microtype}
\microtypecontext{spacing=nonfrench}

% Line numbers
\usepackage[right]{lineno}

% Text layout - adjust as needed
\raggedright
\setlength{\parindent}{0.5cm}
\textwidth 5.25in 
\textheight 8.75in

% Set double spacing
\usepackage{setspace} 
\doublespacing

% Adjust width for specific content
\usepackage{changepage}

% Adjust caption style
\usepackage[aboveskip=1pt,labelfont=bf,labelsep=period,singlelinecheck=off]{caption}

% Remove brackets from references
\makeatletter
\renewcommand{\@biblabel}[1]{\quad#1.}
\makeatother

% Header, footer, and page numbers
\usepackage{lastpage,fancyhdr}
\pagestyle{fancy}
\fancyhf{}
\fancyhead[L]{Preprint - PrimeAI Enhanced Template}
\fancyfoot[C]{\scriptsize Multiplicity Theory © 2025 Ryan O. Van Gelder - Citizen Gardens \\ Licensed Under MIT and CC BY-NC-SA 4.0.}
\fancyfoot[R]{Page \thepage\ of \pageref{LastPage}}
\renewcommand{\footrule}{\hrule height 2pt \vspace{2mm}}

\begin{document}
\title{Quantum Micro-Macro Superposition: \\ Bridging Quantum and Classical Dynamics through Multiplicity Theory}

\author{Ryan O. Van Gelder}
\affil{Citizen Gardens - The Foundation of Multiplicity}
\date{\today}
\maketitle

\begin{abstract}
This paper explores the theoretical and practical implications of achieving quantum micro-macro superposition through Multiplicity Theory. By integrating prime-based encoding, tensor networks, and recursive feedback mechanisms, this work aims to establish a framework for coherent quantum behaviors in macroscopic systems. Applications in quantum computing, life sciences, and hybrid quantum-classical paradigms are examined, along with experimental pathways for validation.
\end{abstract}

\section{Introduction}
Multiplicity Theory offers a robust mathematical foundation for understanding interconnected systems across scales \cite{van2024multiplicity}. Its principles of prime-based encoding, tensor dynamics, and recursive feedback enable novel approaches to bridging quantum and classical paradigms. Here, we investigate whether macroscopic systems, such as biological cells or large molecules, can exhibit quantum behaviors if mediated by Multiplicity Theory.

\section{Mathematical Framework}
The following sections outline the mathematical constructs used in this study, including dynamic multiplicity equations and tensor-based models.

\subsection{Dynamic Multiplicity Equation}
The evolution of macroscopic quantum states is governed by the enhanced Dynamic Multiplicity Equation:
\begin{equation}
\frac{\partial \rho_k}{\partial t} = \alpha_k(t) \rho_k + \beta_k(t) I_k + \gamma_k(t) \sum_j T_{kj} \rho_j + \lambda(t)(\Omega_B(\rho) + \Omega_{FS}(\rho)) + \eta_k \rho_k^2 + \xi_k(t),
\label{eq:dynamic_eq}
\end{equation}
where:
\begin{itemize}
    \item $\rho_k$ represents the state of the $k$-th system component.
    \item $\alpha_k(t)$, $\beta_k(t)$, and $\gamma_k(t)$ are time-dependent parameters.
    \item $T_{kj}$ denotes the tensor coupling matrix.
    \item $\Omega_B(\rho)$ and $\Omega_{FS}(\rho)$ are geometric feedback terms.
    \item $\eta_k$ accounts for nonlinear saturation effects.
    \item $\xi_k(t)$ models stochastic environmental noise.
\end{itemize}

\subsection{Prime-Based Encoding}
Prime-based encoding enables unique and dynamic representations of system states:\newline
\begin{equation}
\phi_k = e^{2\pi i n_k/p_k} \cdot e^{i\theta_k},
\end{equation}
where $n_k$ and $p_k$ are integers representing the system state, and $\theta_k$ encodes phase information \cite{van2024prime}.

\subsection{Tensor Networks for Multi-Scale Modeling}
Tensor networks capture interactions across scales:
\begin{equation}
T(t) = \sum_{i,j,k} T_{ijk} \otimes \phi(p_{ijk}),
\end{equation}
where $\phi(p_{ijk})$ uses prime-based encoding to track state dependencies.

\section{Applications}

\subsection{Quantum Computing}
The proposed framework enhances quantum computing by integrating large-scale coherence:\newline
\begin{equation}
\psi(t) = \sum_{i=1}^N \lambda_i \mu_i \cdot e^{i\theta_i(t)} \cdot v_i,
\end{equation}
where $\lambda_i$ and $\mu_i$ are eigenvalues and multiplicities, respectively, and $v_i$ are eigenvectors \cite{van2024unified}.

\subsection{Quantum Biology}
Biological systems, such as photosynthetic complexes, may leverage quantum coherence \cite{galioto2024dynamic}:
\begin{equation}
M(t) = \sum_{i,j} C_{ij}(t) \cdot \phi_i(t) \phi_j^*(t),
\end{equation}
where $C_{ij}(t)$ is the correlation matrix representing entanglement.

\section{Experimental Validation}

\subsection{Simulation of Macroscopic Quantum States}
Simulations can test the feasibility of macroscopic coherence:\newline
\begin{equation}
P(S, t) = \prod_{i=1}^N \left(p(x_i, t) + \epsilon_i(t)\right),
\end{equation}
where $p(x_i, t)$ represents encoded states, and $\epsilon_i(t)$ models noise.

\subsection{Biological Experiments}
Experimental setups for biological systems include: 
\begin{itemize}
    \item Quantum coherence measurement in photosynthetic proteins using femtosecond spectroscopy.
    \item Observing decoherence suppression in neural circuits under specific conditions.
\end{itemize}

\section{Future Directions}
This framework opens possibilities for:
\begin{itemize}
    \item Large-scale quantum simulations.
    \item Hybrid quantum-classical computing systems.
    \item Exploration of quantum effects in macroscopic biological systems.
\end{itemize}

\nocite{*}
\bibliographystyle{plain}
\bibliography{references}

\end{document}